\section*{Chapitre \ref{ch:g11:subduction-zones} --- Les zones de subduction}

\begin{solution}{exo:g11:subduction-zones:1}
La \omterm{def:g11:subduction-zones:subduction}{subduction} est l'enfoncement d'une plaque de \omterm{def:g11:ocean-ridges:oceanic-lithosphere}{lithosphère océanique} sous
une autre plaque, qui la ramène dans le manteau~; une \omterm{def:g11:subduction-zones:subduction}{zone de subduction} est
une limite de plaques convergente où elle se produit, marquée par une fosse.
C'est la \omterm{def:g11:ocean-ridges:oceanic-lithosphere}{lithosphère océanique} qui s'enfonce~: une fois ancienne, elle est
froide, épaisse et plus dense que l'\omterm{def:g11:lithosphere-asthenosphere:lithosphere}{asthénosphère} située dessous. La croûte
continentale est trop légère pour s'enfoncer.
\end{solution}

\begin{solution}{exo:g11:subduction-zones:2}
(1) Par exemple la région des Tonga et des Fidji, le Pacifique Ouest du
Japon aux Mariannes, l'Indonésie et l'Amérique du Sud. (2) Toujours du côté
de la plaque chevauchante, jamais du côté océanique. (3) Faible~: tous les
séismes de la dorsale médio-atlantique figurant sur la carte se produisent
à moins de \qty{70}{km} de profondeur.
\end{solution}

\begin{solution}{exo:g11:subduction-zones:3}
$7031 \times 0.21 \approx 1480$ séismes à plus de \qty{70}{km} de
profondeur~; $7031 \times 0.069 \approx 485$ à plus de \qty{300}{km}.
\end{solution}

\begin{solution}{exo:g11:subduction-zones:4}
Le \omterm{def:g11:subduction-zones:eclogite}{schiste bleu} est une roche métamorphique issue du basalte de la \omterm{def:g11:ocean-ridges:oceanic-lithosphere}{croûte
océanique} enfoui à haute pression tout en restant relativement froid, et
coloré en bleu par la glaucophane~; l'\omterm{def:g11:subduction-zones:eclogite}{éclogite} est la roche métamorphique
dense que deviennent le basalte et le \omterm{def:g11:ocean-ridges:gabbro}{gabbro} aux profondeurs du manteau,
faite de grenat rouge et de pyroxène vert. Tous deux viennent de la \omterm{def:g11:ocean-ridges:oceanic-lithosphere}{croûte
océanique} et en gardent la composition chimique~; ce qui diffère, ce sont
les minéraux, formés à l'état solide sous les nouvelles conditions de
pression et de température, et la masse volumique.
\end{solution}

\begin{solution}{exo:g11:subduction-zones:5}
Fond océanique~: une \omterm{def:g11:lithosphere-asthenosphere:lithosphere}{lithosphère} ancienne qui avance vers la fosse. Fosse~:
la plaque se courbe vers le bas et commence à s'enfoncer. Avant-arc~: les
sédiments raclés sur la plaque qui s'enfonce s'entassent~; dessous, les deux
plaques glissent l'une contre l'autre sur la limite de plaques, où se
produisent de grands séismes. Arc volcanique~: un alignement de volcans
au-dessus du \omterm{def:g11:subduction-zones:slab}{panneau plongeant}. Arrière-arc~: derrière certains arcs, la
plaque chevauchante est étirée.
\end{solution}

\begin{solution}{exo:g11:subduction-zones:6}
$\tan\delta = (508 - 135)/(800 - 300) = 373/500 \approx 0.746$, donc
$\delta \approx \ang{37}$.
\end{solution}

\begin{solution}{exo:g11:subduction-zones:7}
(1) $\tan\ang{45} = 1$~: à \qty{300}{km} de profondeur. (2) À \qty{670}{km}
de la fosse. (3) $670\sqrt{2} \approx \qty{948}{km}$ de \omterm{def:g11:subduction-zones:slab}{panneau plongeant}.
\end{solution}

\begin{solution}{exo:g11:subduction-zones:8}
(1) $(74.6 + 70.0 + 65.5 + 63.1 + 52.6)/5 \approx \qty{65.2}{mm}$ par an~;
$(47.6 + 31.0 + 29.2 + 22.8 + 19.8 + 15.2 + 14.9 + 10.6)/8 \approx
\qty{23.9}{mm}$ par an~; rapport $65.2/23.9 \approx 2.7$. (2) Les plaques se
déplacent le plus vite quand un \omterm{def:g11:subduction-zones:slab}{panneau plongeant} pend de leur bord~: la
\omterm{def:g11:subduction-zones:slab-pull}{traction du panneau plongeant} est la force la plus intense qui met les
plaques en mouvement, comme l'ont conclu Forsyth et Uyeda.
\end{solution}

\begin{solution}{exo:g11:subduction-zones:9}
(1) \qty{83}{mm} font \qty{0.083}{m}~: $50/0.083 \approx 600$ ans et
$60/0.083 \approx 720$ ans de mouvement de plaque. (2) De 869 à 2011, il
s'est écoulé 1142 ans~: c'est du même ordre. Si la limite de plaques reste
bloquée entre deux grands séismes, le glissement de 2011 a libéré plusieurs
siècles de mouvement accumulé.
\end{solution}

\begin{solution}{exo:g11:subduction-zones:10}
(1) Japon~: $6400 - 3200\,e^{-133/62.8} \approx \qty{6015}{m}$~; Tonga~:
$6400 - 3200\,e^{-99/62.8} \approx \qty{5740}{m}$~; Andes~:
$2500 + 350\sqrt{49} = \qty{4950}{m}$. (2) Le fond qui atteint le Japon~:
c'est le plus ancien, donc la \omterm{def:g11:lithosphere-asthenosphere:lithosphere}{lithosphère} la plus froide et la plus dense,
et le fond le plus profond.
\end{solution}

\begin{solution}{exo:g11:subduction-zones:11}
(1) $3380 \times 1.06 \approx \qty{3583}{kg/m^3}$. (2)
$(3583 - 2900)/2900 \approx 0.24$~: une hausse d'environ \qty{24}{\percent}.
(3) Avant, la croûte est plus légère que le manteau de
$3380 - 2900 = \qty{480}{kg/m^3}$~; après, elle est plus dense d'environ
\qty{203}{kg/m^3}. Elle ne retient plus le \omterm{def:g11:subduction-zones:slab}{panneau plongeant}~: elle
l'entraîne vers le bas.
\end{solution}

\begin{solution}{exo:g11:subduction-zones:12}
(1) $t = 6 \times 10^6 \times 3.16 \times 10^7 \approx 1.9 \times 10^{14}$~s~;
$\sqrt{10^{-6} \times 1.9 \times 10^{14}} \approx \qty{1.4e4}{m}$, environ
\qty{14}{km}. (2) Bien moins que \qty{100}{km}~: en 6 millions d'années,
seule la pellicule externe du \omterm{def:g11:subduction-zones:slab}{panneau plongeant} se réchauffe~; son intérieur
reste froid. (3) La roche froide est rigide et peut se rompre, alors que le
manteau chaud qui l'entoure s'écoule~: les séismes profonds se produisent à
l'intérieur du \omterm{def:g11:subduction-zones:slab}{panneau plongeant} froid.
\end{solution}

\begin{solution}{exo:g11:subduction-zones:13}
(1) $\sqrt{1017^2 + 670^2} \approx \qty{1218}{km}$. (2) \qty{64}{mm} par an
font \qty{64}{km} par million d'années~: $1218/64 \approx 19$ millions
d'années. (3) Environ trois fois plus longtemps que sous les Tonga~: le
\omterm{def:g11:subduction-zones:slab}{panneau plongeant} situé sous les Andes s'enfonce plus lentement et avec un
pendage plus faible.
\end{solution}

\begin{solution}{exo:g11:subduction-zones:14}
(1) Une tranche large de \qty{1}{m} le long de la fosse a un volume de
$10^5 \times 10^6 \times 1 = \qty{1e11}{m^3}$~; son excès de masse vaut
$36 \times 10^{11} = \qty{3.6e12}{kg}$. (2) Son excès de poids vaut
$3.6 \times 10^{12} \times 9.8 \approx \qty{3.5e13}{N}$ pour chaque mètre de
fosse. (3) Le \omterm{def:g11:subduction-zones:slab}{panneau plongeant} est attaché au reste de la plaque, qui est
rigide~: la force se transmet à toute la plaque, qui est entraînée vers la
fosse, d'autant plus que les fosses s'étendent sur des milliers de
kilomètres.
\end{solution}

\begin{solution}{exo:g11:subduction-zones:15}
(1) $(23.6 - 16.5) \times 111 \approx \qty{790}{km}$. (2) La vitesse augmente
de $237 - 104 = \qty{133}{mm}$ par an sur \qty{790}{km}, soit environ
\qty{17}{mm} par an par centaine de kilomètres. (3) La convergence est la
vitesse de la plaque Pacifique moins celle de l'arc. Puisque la vitesse de
la plaque Pacifique est à peu près la même aux deux points, l'arc lui-même
doit se déplacer vers la plaque Pacifique, de plus en plus vite vers le
nord~: l'arrière-arc situé derrière lui est étiré, davantage au nord qu'au
sud.
\end{solution}

\begin{solution}{pb:g11:subduction-zones:1}
\textbf{1.} Vers \qty{671}{km}.

\textbf{2.} Du côté océanique, la plaque ne s'est pas encore enfoncée~: la
\omterm{def:g11:lithosphere-asthenosphere:lithosphere}{lithosphère} y est proche de la surface. Les séismes profonds se produisent
à l'intérieur du \omterm{def:g11:subduction-zones:slab}{panneau plongeant}, qui ne se trouve que sous la plaque
chevauchante.

\textbf{3.} $z(300) = -24 + 273 = \qty{249}{km}$ et
$z(700) = -24 + 637 = \qty{613}{km}$~; $\tan\delta = 364/400 = 0.91$, donc
$\delta \approx \ang{42}$.

\textbf{4.} $x = (100 + 24)/0.91 \approx \qty{136}{km}$ et
$x = (400 + 24)/0.91 \approx \qty{466}{km}$.

\textbf{5.} Une droite de ce pendage partant de la fosse passerait par
$z = 0$ en $x = 0$~; celle-ci passe au-dessus. Le \omterm{def:g11:subduction-zones:slab}{panneau plongeant} n'est
pas rectiligne~: près de la fosse, son pendage est plus faible, et il se
redresse plus bas.

\textbf{6.} \qty{16.8}{cm} par an, \qty{168}{km} par million d'années.

\textbf{7.} $168/69 \approx 2.4$ fois plus rapide.

\textbf{8.} La convergence de part et d'autre de la fosse est la vitesse de
la plaque Pacifique moins la vitesse de l'arc situé de l'autre côté. Si
l'arc se déplace vers la plaque Pacifique, comme il le fait tant que
l'arrière-arc situé derrière lui est étiré, les deux vitesses s'ajoutent, et
la convergence est plus grande que la vitesse de la seule plaque Pacifique.

\textbf{9.} $6400 - 3200\,e^{-99/62.8} \approx \qty{5740}{m}$.

\textbf{10.} Au bout d'environ 60 millions d'années, la \omterm{def:g11:ocean-ridges:oceanic-lithosphere}{lithosphère
océanique} est, dans son ensemble, plus dense que l'\omterm{def:g11:lithosphere-asthenosphere:lithosphere}{asthénosphère} située
dessous~; à 99 millions d'années, elle l'est plus encore. Dès que son bord
descend, son poids entraîne le reste.

\textbf{11.} $(670 + 24)/0.91 \approx \qty{763}{km}$.

\textbf{12.} $\sqrt{763^2 + 670^2} \approx \qty{1015}{km}$.

\textbf{13.} $1015/168 \approx 6.0$ millions d'années.

\textbf{14.} $1015/69 \approx 14.7$ millions d'années.

\textbf{15.} $3380 \times 1.06 \approx \qty{3583}{kg/m^3}$~: la croûte,
autrefois plus légère que le manteau, devient plus dense que lui, et tout le
\omterm{def:g11:subduction-zones:slab}{panneau plongeant} tire plus fort.

\textbf{16.} $t \approx 6.0 \times 10^6 \times 3.16 \times 10^7 \approx
1.9 \times 10^{14}$~s, et $\sqrt{\kappa t} \approx \qty{14}{km}$.

\textbf{17.} \qty{14}{km}, c'est peu à côté de \qty{100}{km}~: seule la
partie externe du \omterm{def:g11:subduction-zones:slab}{panneau plongeant} s'est réchauffée~; son intérieur est
encore froid quand il atteint \qty{670}{km}.

\textbf{18.} Le \omterm{def:g11:subduction-zones:slab}{panneau plongeant} froid reste assez rigide pour se rompre,
jusqu'à son extrémité, alors que le manteau chaud qui l'entoure s'écoule~:
les séismes situés sous les Tonga se produisent à l'intérieur du \omterm{def:g11:subduction-zones:slab}{panneau
plongeant}, jusque vers \qty{671}{km}.

\textbf{19.} Les ondes sismiques révèlent un changement brusque dans le
manteau vers \qty{670}{km}~: leurs vitesses y font un saut.

\textbf{20.} Environ \textbf{6 millions d'années} de la fosse jusqu'à
\qty{670}{km}.
\end{solution}
